% \subsection{CommandHandler}
% \begin{algorithm}[H]
% \caption{CommandHandler (\emph{Interface})}
% \begin{algorithmic}[1]
% \Procedure{Handle}{command : Command\textlangle REQ\textrangle}
%     \State \textbf{Input:} command of type \textlangle REQ\textrangle\
%     \State \textbf{Output:} result of type \textlangle RES\textrangle\
%     \State \textbf{Returns:} the processed result of the command
% \EndProcedure

% \Procedure{Matches}{command : Command\textlangle ?\textrangle}
%     \State \textbf{Input:} command of any type
%     \State \textbf{Output:} Boolean
%     \State handlerType $\gets$ TypeToken of this handler
%     \State payloadType $\gets$ class of command.payload
%     \State \textbf{Return:} true if handlerType is assignable from payloadType, else false
% \EndProcedure
% \end{algorithmic}
% \end{algorithm}

\subsection{CommandHandler}

The \textbf{CommandHandler} interface defines the contract for processing commands of type \textlangle REQ\textrangle\ and producing results of type \textlangle RES\textrangle\ . It serves as the central point of business logic execution in the EOIW pipeline. Each handler encapsulates domain-specific behavior and ensures type-safe command resolution through the \texttt{matches()} function, which dynamically determines compatibility based on generic type inspection. 

Formally, a \emph{CommandHandler} can be expressed as a mapping:
\[
\mathcal{H}: \mathcal{C}_{REQ} \rightarrow \mathcal{R}_{RES}
\]
where $\mathcal{C}_{REQ}$ represents the domain of incoming commands and $\mathcal{R}_{RES}$ represents the resulting output space.  
To enable flexible routing and dynamic dispatch, the handler type $\tau_H$ is resolved at runtime via the \texttt{matches()} method such that:
\[
\tau_H \sqsubseteq \tau_P \implies \text{handler is eligible for command payload type } \tau_P
\]
This design allows the framework to automatically associate commands with the appropriate handlers while maintaining strong compile-time type guarantees.

\begin{algorithm}[H]
\caption{CommandHandler (\emph{Interface})}
\begin{algorithmic}[1]
\Procedure{Handle}{command : Command\textlangle REQ\textrangle}
    \State \textbf{Input:} command of type \textlangle REQ\textrangle
    \State \textbf{Output:} result of type \textlangle RES\textrangle
    \State \textbf{Returns:} the processed result of the command
\EndProcedure

\Procedure{Matches}{command : Command\textlangle ?\textrangle}
    \State \textbf{Input:} command of any type
    \State \textbf{Output:} Boolean
    \State handlerType $\gets$ TypeToken of this handler
    \State payloadType $\gets$ class of command.payload
    \State \textbf{Return:} true if handlerType is assignable from payloadType, else false
\EndProcedure
\end{algorithmic}
\end{algorithm}